% Part: lambda-calculus
% Chapter: church-rosser
% Section: parallel-beta-eta-reduction

\documentclass[../../../include/open-logic-section]{subfiles}

\begin{document}

\olfileid{lam}{cr}{pbe}

\olsection{Reduksi-$\beta\eta$ paralel}

Ing bagean iki kita mbuktekake
sipat Church--Rosser kanggo
reduksi-$\beta\eta$ paralel,
yaiku konsep reduksi paralel
kang cocog karo
reduksi-$\beta\eta$.

\begin{defn}[Reduksi-$\beta\eta$ paralel, $\beredpar$] \ollabel{defn:beredpar}
  \emph{Reduksi-$\beta\eta$ paralel}
  ($\beredpar$) ing suku
  ditetepake kanthi induktif
  kaya mangkene:
  \begin{enumerate}
    \item \ollabel{defn:beredpar1} $x \beredpar x$.
    \item \ollabel{defn:beredpar2} Yen $N \beredpar N'$, mula $\lambd[x][N] \beredpar
      \lambd[x][N']$.
    \item \ollabel{defn:beredpar3} Yen $P \beredpar P'$ lan $Q \beredpar Q'$, mula $PQ \beredpar
      P'Q'$.
    \item \ollabel{defn:beredpar4} Yen $N \beredpar N'$ lan $Q \beredpar Q'$, mula
      $(\lambd[x][N])Q \beredpar \Subst{N'}{Q'}{x}$.
    \item \ollabel{defn:beredpar5} Yen $N \beredpar N'$, mula $\lambd[x][Nx]
      \beredpar N'$, kanthi syarat $x \notin FV(N)$.
  \end{enumerate}
\end{defn}

\begin{thm}\ollabel{thm:refl}
  $M \beredpar M$.
\end{thm}

\begin{proof}
  Latihan.
\end{proof}

\begin{prob}
  Buktekna \olref[lam][cr][pbe]{thm:refl}.
\end{prob}

\begin{defn}[pangembangan lengkap-$\beta\eta$]\ollabel{defn:becd}
  \emph{Pangembangan lengkap-$\beta\eta$}~$\becd{M}$
  saka~$M$ ditetepake
  kaya mangkene:
  \begin{align}
    \becd{x} &= x \ollabel{defn:becd1} \\
    \becd{(\lambd[x][N])} &= \lambd[x][\becd{N}] \ollabel{defn:becd2}\\
    \becd{(PQ)} &= \becd{P}\becd{Q} && \text{yen $P$ dudu abstraksi-$\lambd$}
    \ollabel{defn:becd3} \\
    \becd{((\lambd[x][N])Q)} &= \Subst{\becd{N}}{\becd{Q}}{x}
                             \ollabel{defn:becd4} \\
    \becd{(\lambd[x][Nx])} &= \becd{N} \ollabel{defn:becd5} & \text{yen $x
                                                      \notin FV(N)$}
  \end{align}
  Yen pola abstraksi umum lan
  pola redex eta padha cocog,
  aturan redex eta pungkasan
  kang dienggo.
\end{defn}

\begin{lem}\ollabel{lem:comp}
  Yen $M \beredpar M'$
  lan $R \beredpar R'$,
  mula $\Subst{M}{R}{y}
  \beredpar \Subst{M'}{R'}{y}$.
\end{lem}

\begin{proof}
  Gunakake induksi ing !!{derivation}
  saka $M \beredpar M'$.

  Patang kasus kapisan padha
  polahe karo kasus ing
  \olref[pb]{lem:comp}.
  Yen aturan pungkasan
  \olref{defn:beredpar5},
  $M$ awujud $\lambd[x][Nx]$
  lan $M'$ awujud $N'$
  kanggo sawetara $x$
  lan~$N'$, kanthi
  $x \notin FV(N)$
  lan $N \beredpar N'$.
  Kita kudu nuduhake
  $\Subst{(\lambd[x][Nx])}{R}{y}
  \beredpar \Subst{N'}{R'}{y}$,
  yaiku
  $\lambd[x][\Subst{N}{R}{y} x]
  \beredpar \Subst{N'}{R'}{y}$.
  Iki nggunakake
  \olref{defn:beredpar},
  mligine \olref{defn:beredpar5},
  lan hipotesis induksi.
  Binder kudu dipilih seger
  marang suku pangganti lan
  variabel substitusi supaya
  syarat aturan redex eta
  tetep kacukupan. Justifikasi
  iki durung diwenehake
  ing teks sumber.
\end{proof}

\begin{lem}\ollabel{lem:cont}
  Yen $M \beredpar M'$,
  mula $M' \beredpar \becd{M}$.
\end{lem}

\begin{proof}
  Gunakake induksi ing !!{derivation}
  saka $M \beredpar M'$.

  Patang kasus kapisan
  padha polahe karo kasus ing
  \olref[pb]{lem:cont}.
  Yen aturan pungkasan
  \olref{defn:beredpar5},
  $M$ awujud $\lambd[x][Nx]$
  lan $M'$ awujud $N'$
  kanggo sawetara $x$,
  $N$, $N'$ kanthi
  $x \notin FV(N)$
  lan $N \beredpar N'$.
  Kita kudu nuduhake
  $N' \beredpar
  \becd{(\lambd[x][Nx])}$,
  yaiku $N' \beredpar
  \becd{N}$; iki
  langsung saka hipotesis
  induksi.
  Nanging kasus abstraksi
  kang uga redex eta
  butuh pamriksa dhewe
  amarga loro aturan
  pangembangan lengkap
  bisa tumpang tindih.
\end{proof}

\begin{thm}\ollabel{thm:cr}
  $\beredpar$ nduweni
  sipat Church--Rosser.
\end{thm}

\begin{proof}
  Langsung saka \olref{lem:cont},
  kanthi gumantung marang
  pamriksa kasus-kasus
  kang kasebut ing ndhuwur.
\end{proof}

\end{document}
